% Section 10.1, from lectures/L10/appendix/a_ascii_and_16bit.md.
%
% One correction, made in the lecture too: 10.1.2 said that AVR has no instruction that adds a
% constant, which 10.1.8 then contradicts with adiw; it now says "to a single register".
% The long comment in the 16-bit addition drawing of 10.1.7 is moved to a line of its own, so the
% drawing fits the page.

\section{ASCII-koder och 16-bitarstal}
\label{c10:app:a}

\subsection{Tecken är också tal}
\label{c10:a1}
En mikrodator känner bara till tal. Ett tecken, en bokstav eller en siffra på en display, är ett tal
som alla har kommit överens om att tolka som ett visst tecken. Den överenskommelse som nästan all
elektronik använder heter \textbf{ASCII}, och den gicks igenom i \secref{c1:app:b}. Här behövs
bara en liten del av tabellen:

\begin{narrowtable}{@{}cc@{\hspace{3em}}cc@{}}
  \thead{Tecken} & \thead{ASCII-kod} & \thead{Tecken} & \thead{ASCII-kod} \\
  \midrule
  \code{'0'} & \code{0x30} & \code{'A'} & \code{0x41} \\
  \code{'1'} & \code{0x31} & \code{'B'} & \code{0x42} \\
  \code{'2'} & \code{0x32} & \code{'C'} & \code{0x43} \\
  \dots & \dots & \code{'D'} & \code{0x44} \\
  \code{'8'} & \code{0x38} & \code{'E'} & \code{0x45} \\
  \code{'9'} & \code{0x39} & \code{'F'} & \code{0x46} \\
\end{narrowtable}

\noindent Två saker i tabellen gör resten av det här avsnittet möjligt:
\begin{itemize}
  \item \textbf{Siffrorna ligger i ordning.} Tecknet för siffran \code{n} har koden
    \code{0x30 + n}. Tecknet \code{'7'} är alltså \code{0x37}, och siffran 7 blir ett tecken genom
    att man adderar \code{0x30}.
  \item \textbf{Bokstäverna ligger också i ordning, men inte direkt efter siffrorna.} Efter
    \code{'9'} (\code{0x39}) kommer sju andra tecken (\code{:}, \code{;}, \code{<} och så vidare)
    innan \code{'A'} (\code{0x41}). Den luckan på sju kommer att dyka upp i koden.
\end{itemize}

Assemblern känner till tecken. \code{ldi r16, 'A'} betyder precis samma sak som
\code{ldi r16, 0x41}, och det är ofta det tydligaste sättet att skriva en konstant som är tänkt som
ett tecken.

Varför behövs det här? Så fort ett tal ska visas för en människa, på en display, i ett
terminalfönster eller i en logg, måste det göras om till tecken. Talet \code{0x3C} i ett register
är en byte; texten \enquote{3C} på en skärm är två byte, \code{0x33} och \code{0x43}.

\subsection{En siffra blir ett tecken}
\label{c10:a2}
Ett värde 0--15 ska bli tecknet för sin hexadecimala siffra, \code{'0'}--\code{'9'} eller
\code{'A'}--\code{'F'}.

\textbf{Värdet 0--9:} addera \code{0x30}. Värdet 7 blir \code{0x37}, som är \code{'7'}.

\textbf{Värdet 10--15:} addera \code{0x30} och sju till. Värdet 10 blir
\code{0x0A + 0x07 + 0x30 = 0x41}, som är \code{'A'}. De sju extra är luckan mellan \code{'9'} och
\code{'A'} från \secref{c10:a1}.

AVR saknar en instruktion som adderar en konstant till ett enskilt register. Men att subtrahera ett
negativt tal är samma sak som att addera, så \code{subi r24, -0x30} adderar \code{0x30} (se
\secref{c8:app:a}). Hela omvandlingen blir fyra instruktioner och en etikett:

\needspace{7\baselineskip}
\begin{asmcode}
    cpi r24, 10                     ; A digit 0-9, or a letter A-F?
    brlo digit                      ; 0-9: skip the extra step.
    subi r24, -7                    ; A-F: add 7 more, so that 10 lands on 'A'.
digit:
    subi r24, -0x30                 ; Add 0x30, the ASCII code for '0'.
\end{asmcode}

\noindent Följ två värden genom koden:

\begin{booktable}{@{}p{4.4em}p{6.2em}p{6.6em}LL@{}}
  \thead{\code{r24} före} & \thead{\code{cpi r24, 10}} & \thead{Hopp?} &
    \thead{Efter \code{subi r24, -7}} & \thead{Efter \code{subi r24, -0x30}} \\
  \midrule
  \code{0x03} & 3 < 10, C = 1 & ja, \code{brlo} hoppar & (körs inte) & \code{0x33} = \code{'3'} \\
  \code{0x0C} & 12 ≥ 10, C = 0 & nej & \code{0x13} & \code{0x43} = \code{'C'} \\
\end{booktable}

\begin{aside}
\textbf{Flaggorna efter \code{subi} med ett negativt tal} beter sig som vid en subtraktion, inte
som vid en addition: \code{C} blir 1 när det \emph{inte} blev någon minnessiffra i additionen. Här
spelar det ingen roll, eftersom ingenting läser flaggorna efteråt. Men bygg aldrig ett villkorligt
hopp på flaggorna efter ett sådant \code{subi}.
\end{aside}

\subsection{En byte blir två tecken}
\label{c10:a3}
En byte är två hexadecimala siffror, och varje siffra är fyra bitar, en \textbf{nibble}
(\secref{c1:app:a}). För att skriva ut byten \code{0x3C} ska den höga nibblen \code{3} och den låga
nibblen \code{C} göras om var för sig.

\bookfigure{l10_nibbles}{Byten \code{0x3C} delas i sin höga och sin låga halva, som blir var sitt
tecken.}{c10:fig:nibbles}

\textbf{Den låga nibblen} får man genom att nollställa de fyra höga bitarna med en mask:
\code{andi r24, 0x0F} behåller bit 0--3 och nollställer bit 4--7. \code{0x3C} blir \code{0x0C}.

\textbf{Den höga nibblen} får man med en ny instruktion, \textbf{\code{swap}}, som byter plats på
ett registers två halvor: \code{0x3C} blir \code{0xC3}. Därefter gör samma mask som nyss att
\code{0x03} blir kvar. \code{swap} tar en cykel och påverkar inga flaggor.

Hela programmet, \code{hex_to_ascii.asm}, gör om \code{0x3C} i \code{r16} till tecknen \code{'3'} i
\code{r20} och \code{'C'} i \code{r21}:

\begin{asmcode}
main:
    ldi r16, 0x3C                   ; The byte to convert.

    ; The high nibble: 0x3C -> 0x03 -> '3'.
    mov r24, r16                    ; Work on a copy, so r16 keeps its value.
    swap r24                        ; Swap the two nibbles: 0x3C -> 0xC3.
    andi r24, 0x0F                  ; Keep only the low nibble: 0xC3 -> 0x03.
    cpi r24, 10                     ; A digit 0-9, or a letter A-F?
    brlo high_digit                 ; 0-9: skip the extra step.
    subi r24, -7                    ; A-F: add 7 more, so that 10 lands on 'A'.
high_digit:
    subi r24, -0x30                 ; Add 0x30, the ASCII code for '0'.
    mov r20, r24                    ; r20 = '3' = 0x33.

    ; The low nibble: 0x3C -> 0x0C -> 'C'.
    mov r24, r16
    andi r24, 0x0F                  ; Keep only the low nibble: 0x3C -> 0x0C.
    cpi r24, 10
    brlo low_digit
    subi r24, -7
low_digit:
    subi r24, -0x30
    mov r21, r24                    ; r21 = 'C' = 0x43.
\end{asmcode}

Lägg märke till att samma fem rader står två gånger, bara med olika etiketter. Det är ett tecken på
att de borde vara en \textbf{subrutin}, något man skriver en gång och anropar två gånger. Det är
ämnet för \secref{c10:app:c}, och i \secref{c10:c9} skrivs det här programmet om så.

\subsection{Tillbaka från tecken till värde}
\label{c10:a4}
Omvänt: ett tecken \code{'0'}--\code{'9'} eller \code{'A'}--\code{'F'} ska bli sitt värde 0--15.
Det är samma steg i omvänd ordning:
\begin{enumerate}
  \item Subtrahera \code{0x30}. En siffra är nu klar: \code{'7'} (\code{0x37}) blir 7.
  \item En bokstav blev 17 eller mer: \code{'A'} (\code{0x41}) blir \code{0x11} = 17. Subtrahera
    sju till.
\end{enumerate}

\needspace{7\baselineskip}
\begin{asmcode}
    subi r22, 0x30                  ; '0'-'9' -> 0-9, 'A'-'F' -> 17-22.
    cpi r22, 10
    brlo done                       ; A digit: done.
    subi r22, 7                     ; A letter: 17-22 -> 10-15.
done:
\end{asmcode}

\noindent Det är så en mikrodator tolkar ett tal som någon har skrivit in på ett tangentbord: varje
tangent skickar ett tecken, och programmet gör om tecknen till tal.

\subsection{Att lägga text i minnet}
\label{c10:a5}
Hittills har varje värde legat i ett register. Text brukar i stället ligga i
\textbf{dataminnet}, SRAM, en byte per tecken efter varandra (\secref{c7:app:a}). En ny
instruktion skriver ett register till en adress i dataminnet:

\needspace{4\baselineskip}
\begin{asmcode}
    sts 0x0100, r20                 ; Store r20 at data address 0x0100.
    sts 0x0101, r21                 ; ... and r21 right after it.
\end{asmcode}

\noindent\textbf{\code{sts}} (\emph{store direct to data space}) tar en adress och ett register,
och tar två cykler. Motsatsen, \textbf{\code{lds r16, 0x0100}}, läser tillbaka byten till ett
register. SRAM börjar på adress \code{0x0100}, så de adresserna är lediga att använda för egna
data.

Lägg till de två raderna sist i programmet i \secref{c10:a3}, så hamnar tecknen \code{'3'} och
\code{'C'} efter varandra i minnet. I Microchip Studio syns de i fönstret \textbf{Memory} med
\textbf{data IRAM} valt (\secref{studio:4-3} i bilaga~\ref{app:studio}). Till höger om bytena
\code{33 43} visar fönstret samma byte tolkade som ASCII, och där står texten \code{3C}.

\subsection{Tal större än 255}
\label{c10:a6}
Ett register är åtta bitar, och åtta bitar räcker till talen 0--255. Det räcker långt, men inte
till allt: ett räknevärde på 1000, en fördröjning på en sekund, eller en adress i ett minne på
2048 byte.

Lösningen är att använda \textbf{två register tillsammans}, som ett 16-bitars tal. Det ena håller
den \textbf{höga byten} (bit 15--8) och det andra den \textbf{låga byten} (bit 7--0). Ett sådant
\textbf{registerpar} skrivs med det höga registret först, till exempel \code{r25:r24}.

\bookfigure{l10_register_pair}{Registerparet \code{r25:r24} som håller talet
1000.}{c10:fig:register_pair}

Talet 1000 är \code{0x03E8}. Den höga byten är \code{0x03} och den låga \code{0xE8}, och värdet är
\code{0x03 · 256 + 0xE8 = 768 + 232 = 1000}. En etta i den höga byten är alltså värd 256.

Assemblern delar upp ett 16-bitars tal åt dig, med \code{low()} och \code{high()}:

\needspace{4\baselineskip}
\begin{asmcode}
    ldi r24, low(1000)              ; r24 = 0xE8, the low byte.
    ldi r25, high(1000)             ; r25 = 0x03, the high byte.
\end{asmcode}

\noindent Med 16 bitar går det att räkna till \mbox{\code{2^16 - 1 = 65 535}}. Paren \code{r25:r24},
\code{r27:r26}, \code{r29:r28} och \code{r31:r30} används oftast, av skäl som \secref{c10:a8}
visar.

\subsection{Addition och subtraktion med 16 bitar}
\label{c10:a7}
Processorn räknar fortfarande åtta bitar i taget. En 16-bitars addition görs därför i två steg, på
samma sätt som när man adderar för hand med minnessiffra (\secref{c1:a7}): först de låga bytena,
sedan de höga bytena \textbf{plus minnessiffran} från det första steget.

Minnessiffran hamnar i flaggan \code{C}, och det finns en instruktion som tar med den:
\textbf{\code{adc}} (\emph{add with carry}) beräknar \code{Rd = Rd + Rr + C}.

Exempel: 1000 + 2000 = 3000, eller \code{0x03E8 + 0x07D0}.

\needspace{11\baselineskip}
\begin{textdrawing}
                 r25    r24
                0x03   0xE8      1000
             +  0x07   0xD0      2000
             ---------------
  add r24:             0x1B8  -> r24 = 0xB8, C = 1
                                 (0xE8 + 0xD0 = 0x1B8 får inte plats)
  adc r25:     0x03 + 0x07 + 1 = 0x0B
             ---------------
                0x0B   0xB8      3000
\end{textdrawing}

\needspace{8\baselineskip}
\begin{asmcode}
    ldi r24, low(1000)              ; r25:r24 = 1000 = 0x03E8.
    ldi r25, high(1000)
    ldi r22, low(2000)              ; r23:r22 = 2000 = 0x07D0.
    ldi r23, high(2000)
    add r24, r22                    ; Low bytes: 0xE8 + 0xD0 = 0x1B8. r24 = 0xB8, C = 1.
    adc r25, r23                    ; High bytes and the carry: 0x03 + 0x07 + 1 = 0x0B.
\end{asmcode}

\noindent Utan \code{adc}, med \code{add} även för de höga bytena, hade svaret blivit \code{0x0AB8}
= 2744: minnessiffran hade försvunnit, och felet hade varit exakt 256. Ett fel på 256 i ett
16-bitars resultat är nästan alltid en glömd minnessiffra.

Subtraktion fungerar likadant. \textbf{\code{sub}} tar de låga bytena, och \textbf{\code{sbc}}
(\emph{subtract with carry}) tar de höga bytena minus lånet i \code{C}: \code{Rd = Rd - Rr - C}.

\needspace{4\baselineskip}
\begin{asmcode}
    sub r20, r22                    ; Low bytes first; a borrow lands in C ...
    sbc r21, r23                    ; ... and sbc subtracts it from the high bytes.
\end{asmcode}

\noindent\textbf{Ordningen spelar roll:} de låga bytena först, alltid. Det är den operationen som
ger \code{C}, och den höga behöver det.

\subsection{\code{adiw}, \code{sbiw} och \code{movw}}
\label{c10:a8}
Tre instruktioner arbetar på ett helt registerpar på en gång:

\begin{booktable}{@{}p{7.4em}Lc@{}}
  \thead{Instruktion} & \thead{Gör} & \thead{Cykler} \\
  \midrule
  \code{adiw r24, 1} & \code{r25:r24 = r25:r24 + 1}. Konstanten får vara 0--63. & 2 \\
  \code{sbiw r24, 1} & \code{r25:r24 = r25:r24 - 1}. Konstanten får vara 0--63. & 2 \\
  \code{movw r18, r24} & Kopierar paret: \code{r19:r18 = r25:r24}. & 1 \\
\end{booktable}

Man skriver bara det låga registret i paret; det höga följer med. \code{adiw} och \code{sbiw}
fungerar bara på de fyra paren \code{r25:r24}, \code{r27:r26}, \code{r29:r28} och \code{r31:r30},
vilket är varför de paren används mest. \code{movw} fungerar på alla par vars låga register har
jämnt nummer.

Det viktigaste med \code{sbiw} är flaggan \textbf{Z}: den blir 1 när \textbf{hela}
16-bitarstalet blev noll, inte bara den ena byten. Därför kan \code{sbiw} följt av \code{brne}
räkna ned ett 16-bitars tal till noll, vilket är precis vad en lång fördröjning behöver
(\secref{c10:a10}).

Programmet \code{add16.asm} visar alla tre: efter additionen i \secref{c10:a7} kopieras 3000 till
\code{r19:r18} med \code{movw}, och \code{r25:r24} blir sedan 3001 med \code{adiw} och 2999 med
\code{sbiw}.

\subsection{Att jämföra 16-bitarstal}
\label{c10:a9}
En jämförelse görs också byte för byte. \code{cpi} jämför de låga bytena, och \textbf{\code{cpc}}
(\emph{compare with carry}) jämför de höga bytena och tar hänsyn till resultatet av den första
jämförelsen. Efter paret säger \code{Z} om talen var lika, och \code{C} om det första var mindre,
precis som för 8-bitarstal.

\needspace{8\baselineskip}
\begin{asmcode}
    ldi r16, high(1000)             ; cpc compares with a register, not a constant.
count:
    adiw r24, 1                     ; r25:r24 = r25:r24 + 1.
    cpi r24, low(1000)              ; Compare the low bytes ...
    cpc r25, r16                    ; ... then the high bytes.
    brne count                      ; Not 1000 yet: again.
\end{asmcode}

\noindent\code{cpc} finns bara i en form som jämför två register. Det finns ingen \code{cpci}, så
den höga byten av konstanten måste först laddas i ett eget register, här \code{r16}.

Loopen räknar \code{r25:r24} från 0 till 1000. Varje varv tar \code{2 + 1 + 1 + 2 = 6} cykler,
utom det sista, där \code{brne} inte hoppar och tar en cykel: \code{1000 · 6 - 1 = 5999} cykler.

\subsection{En lång tidsfördröjning}
\label{c10:a10}
Fördröjningarna i \secref{c9:app:b} räknade med 8-bitars register. En loop med \code{dec} och
\code{brne} kan gå högst 256 varv, och med 3 cykler per varv blir det ungefär 768 cykler, 48~µs.
Två sådana loopar i varandra räcker till ungefär 200\,000 cykler, drygt 12~ms. För en sekund,
16\,000\,000 cykler vid 16~MHz, behövs mer.

Med en 16-bitars loop kan den inre loopen gå upp till 65\,535 varv. En sådan loop med \code{sbiw}
och \code{brne} tar 4 cykler per varv:

\needspace{7\baselineskip}
\begin{asmcode}
    ldi r24, low(INNER_TURNS)       ; 1 cycle.
    ldi r25, high(INNER_TURNS)      ; 1 cycle.
inner:
    sbiw r24, 1                     ; 2 cycles.
    brne inner                      ; 2 cycles while not 0, 1 the last time.
\end{asmcode}

\noindent Med \code{N} varv tar loopen \code{4N - 1} cykler: \code{N} varv om 4 cykler, minus en
eftersom den sista \code{brne} inte hoppar. Det räcker till ungefär 262\,000 cykler, 16~ms. För en
hel sekund läggs ytterligare en loop utanför, som kör den inre loopen \code{M} gånger. Hela
fördröjningen i \code{one_second_delay.asm}:

\needspace{16\baselineskip}
\begin{asmcode}
.equ OUTER_TURNS = 80               ; Outer loop: 80 turns.
.equ INNER_TURNS = 49999            ; Inner loop: 49 999 turns of 4 cycles each.

delay_begin:
    ldi r20, OUTER_TURNS            ; 1 cycle.
outer:
    ldi r24, low(INNER_TURNS)       ; 1 cycle.
    ldi r25, high(INNER_TURNS)      ; 1 cycle.
inner:
    sbiw r24, 1                     ; 2 cycles: r25:r24 = r25:r24 - 1.
    brne inner                      ; 2 cycles while r25:r24 is not 0, 1 the last time.
    dec r20                         ; 1 cycle.
    brne outer                      ; 2 cycles while r20 is not 0, 1 the last time.
delay_end:
\end{asmcode}

\noindent Räkna rad för rad, som i kapitel~\ref{c9:ch}:

\begin{narrowtable}{@{}lccc@{}}
  \thead{Instruktion} & \thead{Cykler} & \thead{Gånger} & \thead{Summa} \\
  \midrule
  \code{ldi r20} & 1 & 1 & 1 \\
  \code{ldi r24}, \code{ldi r25} & 1 + 1 & \code{M} & \code{2M} \\
  inre loopen & \code{4N - 1} & \code{M} & \code{M(4N - 1)} \\
  \code{dec r20} & 1 & \code{M} & \code{M} \\
  \code{brne outer}, hoppar & 2 & \code{M - 1} & \code{2M - 2} \\
  \code{brne outer}, hoppar inte & 1 & 1 & 1 \\
\end{narrowtable}

\noindent Summan är \code{M(4N + 4)}. Ett varv i den yttre loopen tar alltså \code{4N + 4} cykler,
och med \mbox{\code{N = 49 999}} blir det exakt 200\,000 cykler, 12,5~ms. Då behövs
\mbox{\code{M = 16 000 000}} \mbox{\code{/ 200 000}} \mbox{\code{= 80}} varv för en sekund:

\needspace{3\baselineskip}
\begin{textdrawing}
M(4N + 4) = 80 · (4 · 49 999 + 4) = 80 · 200 000 = 16 000 000 cykler = 1,000000 s
\end{textdrawing}

\noindent Simulatorn mäter också 16\,000\,000 cykler mellan \code{delay_begin} och
\code{delay_end}. Så väljer man konstanterna för en önskad tid:
\begin{enumerate}
  \item Räkna om tiden till cykler: tid · 16\,000\,000.
  \item Välj ett jämnt varv för den yttre loopen, till exempel 200\,000 cykler, och räkna ut
    \code{N} ur \mbox{\code{4N + 4 = 200 000}}.
  \item Dela antalet cykler med varvtiden, så får du \code{M}. \code{M} måste vara 1--255.
\end{enumerate}

\begin{aside}
\textbf{Namnen på konstanter och etiketter får inte krocka.} Assemblern skiljer inte på stora och
små bokstäver, så \code{.equ OUTER = 80} och etiketten \code{outer:} är samma namn, och programmet
assembleras inte. Därför heter konstanterna \code{OUTER_TURNS} och \code{INNER_TURNS}.
\end{aside}
